\begin{thebibliography}{43}

\bibitem[1]{ref1}
Noga Alon, Oded Goldreich, Johan Håstad, and René Peralta. Simple constructions of almost k-wise independent random variables. In 31st Annual Symposium on Foundations of Computer Science, St. Louis, Missouri, USA, October 22-24, 1990, Volume II, pages 544–553. IEEE Computer Society, 1990. doi:10.1109/FSCS.1990.89575.

\bibitem[2]{ref2}
Benny Applebaum and Pavel Raykov. Fast pseudorandom functions based on expander graphs. In Theory of Cryptography - 14th International Conference, TCC 2016-B, Beijing, China, October 31 - November 3, 2016, Proceedings, Part I, volume 9985 of Lecture Notes in Computer Science, pages 27–56, 2016. doi:10.1007/978-3-662-53641-4\_2.

\bibitem[3]{ref3}
Sanjeev Arora and Boaz Barak. Computational Complexity - A Modern Approach. Cambridge University Press, 2009.

\bibitem[4]{ref4}
Abhishek Banerjee, Chris Peikert, and Alon Rosen. Pseudorandom functions and lattices. In Advances in Cryptology - EUROCRYPT 2012 - 31st Annual International Conference on the Theory and Applications of Cryptographic Techniques, Cambridge, UK, April 15-19, 2012. Proceedings, volume 7237 of Lecture Notes in Computer Science, pages 719–737. Springer, 2012. doi:10.1007/978-3-642-29011-4\_42.

\bibitem[5]{ref5}
Andrej Bogdanov and Alon Rosen. Pseudorandom functions: Three decades later. In Tutorials on the Foundations of Cryptography, pages 79–158. Springer International Publishing, 2017. doi:10.1007/978-3-319-57048-8\_3.

\bibitem[6]{ref6}
Larry Carter and Mark N. Wegman. Universal classes of hash functions. J. Comput. Syst. Sci., 18(2):143–154, 1979. doi:10.1016/0022-0000(79)90044-8.

\bibitem[7]{ref7}
L. Sunil Chandran. A high girth graph construction. SIAM J. Discret. Math., 16(3):366–370, 2003. doi:10.1137/S0895480101387893.

\bibitem[8]{ref8}
Lijie Chen, Shuichi Hirahara, Igor Carboni Oliveira, Ján Pich, Ninad Rajgopal, and Rahul Santhanam. Beyond natural proofs: Hardness magnification and locality. In 11th Innovations in Theoretical Computer Science Conference, ITCS 2020, January 12-14, 2020, Seattle, Washington, USA, volume 151 of LIPIcs, pages 70:1–70:48. Schloss Dagstuhl - Leibniz-Zentrum für Informatik, 2020. doi:10.4230/LIPIcs.ITCS.2020.70.

\bibitem[9]{ref9}
Lijie Chen, Ce Jin, and R. Ryan Williams. Hardness magnification for all sparse NP languages. In 60th IEEE Annual Symposium on Foundations of Computer Science, FOCS 2019, Baltimore, Maryland, USA, November 9-12, 2019, pages 1240–1255. IEEE Computer Society, 2019. doi:10.1109/FOCS.2019.00077.

\bibitem[10]{ref10}
Lijie Chen, Ce Jin, and R. Ryan Williams. Sharp threshold results for computational complexity. In Proccedings of the 52nd Annual ACM SIGACT Symposium on Theory of Computing, STOC 2020, Chicago, IL, USA, June 22-26, 2020, pages 1335–1348. ACM, 2020. doi:10.1145/3357713.3384283.

\bibitem[11]{ref11}
Ruiwen Chen and Valentine Kabanets. Correlation bounds and \#SAT algorithms for small linear-size circuits. Theor. Comput. Sci., 654:2–10, 2016. doi:10.1016/j.tcs.2016.05.005.

\bibitem[12]{ref12}
Evgeny Demenkov and Alexander S. Kulikov. An elementary proof of a 3n - o(n) lower bound on the circuit complexity of affine dispersers. In Mathematical Foundations of Computer Science 2011 - 36th International Symposium, MFCS 2011, Warsaw, Poland, August 22-26, 2011. Proceedings, volume 6907 of Lecture Notes in Computer Science, pages 256–265. Springer, 2011. doi:10.1007/978-3-642-22993-0\_25.

\bibitem[13]{ref13}
Zhiyuan Fan, Jiatu Li, and Tianqi Yang. The exact complexity of pseudorandom functions and the black-box natural proof barrier for bootstrapping results in computational complexity. Electron. Colloquium Comput. Complex., page 125, 2021. To appear in STOC 2022. URL: https://eccc.weizmann.ac.il/report/2021/125.

\bibitem[14]{ref14}
Magnus Gausdal Find, Alexander Golovnev, Edward A. Hirsch, and Alexander S. Kulikov. A better-than-3n lower bound for the circuit complexity of an explicit function. In IEEE 57th Annual Symposium on Foundations of Computer Science, FOCS 2016, 9-11 October 2016, Hyatt Regency, New Brunswick, New Jersey, USA, pages 89–98. IEEE Computer Society, 2016. doi:10.1109/FOCS.2016.19.

\bibitem[15]{ref15}
Lance Fortnow. A simple proof of Toda’s theorem. Theory Comput., 5(1):135–140, 2009. doi:10.4086/toc.2009.v005a007.

\bibitem[16]{ref16}
Oded Goldreich, Shafi Goldwasser, and Silvio Micali. How to construct random functions. J. ACM, 33(4):792–807, 1986. doi:10.1145/6490.6503.

\bibitem[17]{ref17}
Alexander Golovnev, Edward A. Hirsch, Alexander Knop, and Alexander S. Kulikov. On the limits of gate elimination. J. Comput. Syst. Sci., 96:107–119, 2018. doi:10.1016/j.jcss. 2018.04.005.

\bibitem[18]{ref18}
Johan Håstad. The shrinkage exponent of de Morgan formulas is 2. SIAM J. Comput., 27(1):48–64, 1998. doi:10.1137/S0097539794261556.

\bibitem[19]{ref19}
Johan Håstad, Russell Impagliazzo, Leonid A. Levin, and Michael Luby. A pseudorandom generator from any one-way function. SIAM J. Comput., 28(4):1364–1396, 1999. doi: 10.1137/S0097539793244708.

\bibitem[20]{ref20}
William Hesse, Eric Allender, and David A. Mix Barrington. Uniform constant-depth threshold circuits for division and iterated multiplication. J. Comput. Syst. Sci., 65(4):695–716, 2002. doi:10.1016/S0022-0000(02)00025-9.

\bibitem[21]{ref21}
Shuichi Hirahara and Rahul Santhanam. On the average-case complexity of MCSP and its variants. In 32nd Computational Complexity Conference, CCC 2017, July 6-9, 2017, Riga, Latvia, volume 79 of LIPIcs, pages 7:1–7:20. Schloss Dagstuhl - Leibniz-Zentrum für Informatik, 2017. doi:10.4230/LIPIcs.CCC.2017.7.

\bibitem[22]{ref22}
Russell Impagliazzo, Raghu Meka, and David Zuckerman. Pseudorandomness from shrinkage. J. ACM, 66(2):11:1–11:16, 2019. doi:10.1145/3230630.

\bibitem[23]{ref23}
Yuval Ishai, Eyal Kushilevitz, Rafail Ostrovsky, and Amit Sahai. Cryptography with constant computational overhead. In Proceedings of the 40th Annual ACM Symposium on Theory of Computing, Victoria, British Columbia, Canada, May 17-20, 2008, pages 433–442. ACM, 2008. doi:10.1145/1374376.1374438.

\bibitem[24]{ref24}
Kazuo Iwama and Hiroki Morizumi. An explicit lower bound of 5n - o(n) for boolean circuits. In Mathematical Foundations of Computer Science 2002, 27th International Symposium, MFCS 2002, Warsaw, Poland, August 26-30, 2002, Proceedings, volume 2420 of Lecture Notes in Computer Science, pages 353–364. Springer, 2002. doi:10.1007/3-540-45687-2\_29.

\bibitem[25]{ref25}
Eyal Kaplan, Moni Naor, and Omer Reingold. Derandomized constructions of k-wise (almost) independent permutations. Algorithmica, 55(1):113–133, 2009. doi:10.1007/ s00453-008-9267-y.

\bibitem[26]{ref26}
Oded Lachish and Ran Raz. Explicit lower bound of 4.5n - o(n) for boolena circuits. In Jeffrey Scott Vitter, Paul G. Spirakis, and Mihalis Yannakakis, editors, Proceedings on 33rd Annual ACM Symposium on Theory of Computing, July 6-8, 2001, Heraklion, Crete, Greece, pages 399–408. ACM, 2001. doi:10.1145/380752.380832.

\bibitem[27]{ref27}
Felix Lazebnik and Vasiliy A. Ustimenko. Explicit construction of graphs with an arbitrary large girth and of large size. Discret. Appl. Math., 60(1-3):275–284, 1995. doi:10.1016/ 0166-218X(94)00058-L.

\bibitem[28]{ref28}
Jiatu Li and Tianqi Yang. 3.1n - o(n) circuit lower bounds for explicit functions. Electron. Colloqu ium Comput. Complex., page 23, 2021. To appear in STOC 2022. URL: https: //eccc.weizmann.ac.il/report/2021/023.

\bibitem[29]{ref29}
Richard J. Lipton and Ryan Williams. Amplifying circuit lower bounds against polynomial time, with applications. Comput. Complex., 22(2):311–343, 2013. doi:10.1007/s00037-013-0069-5.

\bibitem[30]{ref30}
Dylan M. McKay, Cody D. Murray, and R. Ryan Williams. Weak lower bounds on resource- bounded compression imply strong separations of complexity classes. In Proceedings of the 51st Annual ACM SIGACT Symposium on Theory of Computing, STOC 2019, Phoenix, AZ, USA, June 23-26, 2019, pages 1215–1225. ACM, 2019. doi:10.1145/3313276.3316396.

\bibitem[31]{ref31}
Eric Miles and Emanuele Viola. Substitution-permutation networks, pseudorandom functions, and natural proofs. J. ACM, 62(6):46:1–46:29, 2015. doi:10.1145/2792978.

\bibitem[32]{ref32}
Ketan Mulmuley. On P vs. NP and geometric complexity theory: Dedicated to sri ramakrishna. J. ACM, 58(2):5:1–5:26, 2011. doi:10.1145/1944345.1944346.

\bibitem[33]{ref33}
Moni Naor and Omer Reingold. Number-theoretic constructions of efficient pseudo-random functions. J. ACM, 51(2):231–262, 2004. doi:10.1145/972639.972643.

\bibitem[34]{ref34}
Igor Carboni Oliveira, Ján Pich, and Rahul Santhanam. Hardness magnification near state-of- the-art lower bounds. Theory Comput., 17:1–38, 2021.

\bibitem[35]{ref35}
Igor Carboni Oliveira and Rahul Santhanam. Hardness magnification for natural problems. In 59th IEEE Annual Symposium on Foundations of Computer Science, FOCS 2018, pages 65–76, 2018. doi:10.1109/FOCS.2018.00016.

\bibitem[36]{ref36}
Anna Pagh and Rasmus Pagh. Uniform hashing in constant time and optimal space. SIAM J. Comput., 38(1):85–96, 2008. doi:10.1137/060658400.

\bibitem[37]{ref37}
Christos H. Papadimitriou and Stathis Zachos. Two remarks on the power of counting. In Theoretical Computer Science, 6th GI-Conference, Dortmund, Germany, January 5-7, 1983, Proceedings, volume 145 of Lecture Notes in Computer Science, pages 269–276. Springer, 1983. doi:10.1007/BFb0009651.

\bibitem[38]{ref38}
Alexander A. Razborov and Steven Rudich. Natural proofs. J. Comput. Syst. Sci., 55(1):24–35, 1997. doi:10.1006/jcss.1997.1494.

\bibitem[39]{ref39}
Victor Shoup. New algorithms for finding irreducible polynomials over finite fields. In 29th Annual Symposium on Foundations of Computer Science, White Plains, New York, USA, 24-26 October 1988, pages 283–290. IEEE Computer Society, 1988. doi:10.1109/SFCS.1988.21944.

\bibitem[40]{ref40}
Daniel A. Spielman. Linear-time encodable and decodable error-correcting codes. IEEE Trans. Inf. Theory, 42(6):1723–1731, 1996. doi:10.1109/18.556668.

\bibitem[41]{ref41}
Avishay Tal. Shrinkage of de Morgan formulae by spectral techniques. In 55th IEEE Annual Symposium on Foundations of Computer Science, FOCS 2014, Philadelphia, PA, USA, October 18-21, 2014, pages 551–560. IEEE Computer Society, 2014. doi:10.1109/FOCS.2014.65.

\bibitem[42]{ref42}
Leslie G. Valiant and Vijay V. Vazirani. NP is as easy as detecting unique solutions. Theor. Comput. Sci., 47(3):85–93, 1986. doi:10.1016/0304-3975(86)90135-0.

\bibitem[43]{ref43}
Emanuele Viola. The communication complexity of addition. Comb., 35(6):703–747, 2015. doi:10.1007/s00493-014-3078-3.
\end{thebibliography}
